\section*{Bab \ref{ch:b3:adaptive-immunity} --- Imunitas Adaptif dan Vaksinasi}

\begin{solution}{exo:b3:adaptive-immunity:1}
Satu \omterm{def:b3:adaptive-immunity:lymphocytes}{limfosit}, satu kekhasan (Nossal dan Lederberg: tiap sel
menyekresikan \omterm{def:b3:adaptive-immunity:antibody}{antibodi} terhadap satu \omterm{def:b3:adaptive-immunity:lymphocytes}{antigen} saja); reseptornya dibuat sebelum
\omterm{def:b3:adaptive-immunity:lymphocytes}{antigennya} dijumpai (Tonegawa: gennya disusun ulang di dalam sel~B, bukan
sebagai tanggapan terhadap \omterm{def:b3:adaptive-immunity:lymphocytes}{antigen}); \omterm{def:b3:adaptive-immunity:lymphocytes}{antigennya} memilih lalu mengembangkan klona
yang cocok menjadi sel efektor dan \omterm{def:b3:adaptive-immunity:response}{sel ingatan}; klona reaktif diri dihapus selama
perkembangannya.
\end{solution}

\begin{solution}{exo:b3:adaptive-immunity:2}
Dua rantai berat dan dua rantai ringan, masing-masing dengan sebuah \omterm{def:b3:structural-biology:levels}{domain}
berubah-ubah di ujungnya dan \omterm{def:b3:structural-biology:levels}{domain} tetap di belakangnya; dua tapak pengikat
serupa yang dibentuk enam gelung hipervariabel; sebuah batang Fc yang dibaca sel
lain dan \omterm{def:b3:innate-immunity:complement}{komplemen}. IgM: \omterm{def:b3:adaptive-immunity:antibody}{antibodi} pertama, pentamer, pengaktif \omterm{def:b3:innate-immunity:complement}{komplemen}. IgG:
\omterm{def:b3:adaptive-immunity:antibody}{antibodi} utama darah dan jaringan, \omterm{def:b3:innate-immunity:complement}{opsonisasi}, penetralan, menyeberangi plasenta.
IgA: dimer yang disekresikan ke permukaan mukosa dan ke dalam susu. IgE: terikat
pada sel mast, melawan parasit, perantara alergi. IgD: pada sel~B
naif, yang fungsinya kurang dikenal.
\end{solution}

\begin{solution}{exo:b3:adaptive-immunity:3}
Kelas~I: semua sel berinti; peptida dari protein sitosolik yang dipotong
proteasom; 8--10 residu; dibaca sel~T sitotoksik CD8. Kelas~II:
\omterm{def:b3:innate-immunity:innate}{sel dendritik}, \omterm{def:b3:innate-immunity:innate}{makrofag}, sel~B; peptida dari protein yang diambil
lalu dicerna di dalam \omterm{def:b3:membrane-traffic:endocytosis}{endosom}; 13--25 residu; dibaca sel~T penolong CD4.
\end{solution}

\begin{solution}{exo:b3:adaptive-immunity:4}
$R_{0}$ adalah cacah rata-rata kasus yang dihasilkan satu kasus di dalam populasi
yang seluruhnya rentan; ambangnya adalah $p_{c} = 1 - 1/R_{0}$. Gondongan
\qty{80}{\%}; batuk rejan \qty{92}{\%}; influenza 1918 \qty{60}{\%}.
\end{solution}

\begin{solution}{exo:b3:adaptive-immunity:5}
$50\times 20\times 5 = 5000$ rantai berat, $60\times 4 = 240$ rantai ringan:
$1.2\times 10^{6}$ gabungan; dengan \omterm{def:b3:adaptive-immunity:antibody}{keanekaan sambungan} $1.2\times
10^{10}$ --- yakni sepuluh kali cacah sel~B, sehingga hewan itu mengungkapkan
paling banyak sepersepuluh dari apa yang dapat dibuatnya, dan hampir tiap sel
mengusung sebuah reseptor yang unik.
\end{solution}

\begin{solution}{exo:b3:adaptive-immunity:6}
$1000\times 10^{-3} = 1$ mutasi tiap sel tiap pembelahan; $20$ dalam dua puluh
pembelahan. Di antara $10^{5}$ sel, $2\times 10^{6}$ peristiwa mutasi jatuh
pada $3000$ substitusi tunggal yang mungkin: masing-masing tercontoh ratusan
kali, ditambah gabungan rangkap dua dan rangkap tiganya --- \omterm{def:b3:adaptive-immunity:germinal-centre}{pusat germinalnya}
menjelajahi ketetanggaan reseptornya sampai tuntas.
\end{solution}

\begin{solution}{exo:b3:adaptive-immunity:7}
Sel~T terpilih di dalam \omterm{def:b3:adaptive-immunity:tolerance}{timus} untuk mengenali peptida pada alel MHC
inangnya sendiri; alel orang lain punya alur berbentuk lain
yang memegang peptida lain pada arah lain, sehingga
peptida \omterm{def:b3:virology:virus}{virus} yang sama tidak terlihat. Molekul MHC asing sebuah cangkok
dikenali sendiri oleh sebagian besar sel~T sebagai ``MHC diri dengan
sebuah peptida yang aneh'', karena itulah penolakannya. Polimorfismenya bertahan karena
sebuah populasi dengan banyak alel menyajikan contoh yang lebih luas dari peptida
patogen mana pun dan tidak ada patogen yang dapat lolos dari semua orang; heterozigot
menyajikan jangkauan dua kali lipat sehingga lebih dipihak.
\end{solution}

\begin{solution}{exo:b3:adaptive-immunity:8}
(a) Tanpa \omterm{def:b3:adaptive-immunity:antibody}{rekombinasi V(D)J}: tidak ada sel~B atau sel~T, yakni imunodefisiensi
gabungan yang berat. (b) Tanpa AIRE: sel \omterm{def:b3:adaptive-immunity:tolerance}{timusnya} gagal memajang \omterm{def:b3:adaptive-immunity:lymphocytes}{antigen}
jaringan, sel~T reaktif diri lolos, lalu \omterm{def:b3:adaptive-immunity:tolerance}{autoimunitas} terhadap banyak
organ. (c) Tanpa \omterm{def:b3:adaptive-immunity:tolerance}{FoxP3}: tidak ada sel~T pengatur, lalu \omterm{def:b3:adaptive-immunity:tolerance}{autoimunitas} dan alergi
yang mematikan dan berawal dini. (d) Tanpa AID: tidak ada \omterm{def:b3:adaptive-immunity:antibody}{penukaran kelas} atau
hipermutasi --- IgM berlimpah, hampir tidak ada IgG atau IgA, \omterm{def:b3:adaptive-immunity:antibody}{antibodi}
berafinitas rendah, infeksi berulang. (e) Tanpa sel~T CD4: tidak ada pertolongan bagi
sel~B atau \omterm{def:b3:innate-immunity:innate}{makrofag}, tanggapan \omterm{def:b3:adaptive-immunity:antibody}{antibodinya} buruk dan ia mudah terkena
infeksi oportunistik --- yakni imunodefisiensi AIDS.
\end{solution}

\begin{solution}{exo:b3:adaptive-immunity:9}
$p_{c} = 1 - 1/6 = 0.83$; cakupannya $0.83/0.9 = 93\,\%$. Pada cakupan \qty{80}{\%}
bagian yang kebal adalah $0.72$, $R_{\text{eff}} = 6\times 0.28
= 1.7$. Di antara yang rentan, wabahnya menyebar seolah-olah $R_{0}$ adalah
$1.7$ (yang kebal sekadar mengencerkan sentuhannya), dan hubungan ukuran
akhirnya $s_{\infty} - \ln s_{\infty}/1.7 = 1$ memberi $s_{\infty}
\approx 0.31$: sekitar \qty{69}{\%} dari yang rentan, yakni \qty{19}{\%}
populasinya, terinfeksi.
\end{solution}

\begin{solution}{exo:b3:adaptive-immunity:10}
Dari $\mathrm{d}s/\mathrm{d}t = -\beta si$ dan $\mathrm{d}i/\mathrm{d}t
= \beta si - \gamma i$, $\mathrm{d}i/\mathrm{d}s = -1 + \gamma/(\beta s)$,
sehingga $i + s - (\ln s)/R_{0}$ kekal; dengan $s = 1$, $i = 0$ pada
awalnya dan $i = 0$, $s = s_{\infty}$ pada akhirnya, $s_{\infty} - \ln
s_{\infty}/R_{0} = 1$. Penyelesaiannya: $R_{0} = 1.5$, $s_{\infty} = 0.42$
(\qty{58}{\%} terinfeksi); $2$, $0.20$ (\qty{80}{\%}); $3$, $0.06$
(\qty{94}{\%}); $5$, $0.007$ (\qty{99.3}{\%}). Tidak semua orang: ketika
yang rentan terpakai habis, bilangan reproduksi efektifnya jatuh di bawah
satu selagi sebagian masih tersisa, dan wabahnya padam sebelum mencapainya.
\end{solution}

\begin{solution}{exo:b3:adaptive-immunity:11}
$1.5^{n} = 1000$: $n = \ln 1000/\ln 1.5 \approx 17$ perbaikan
berturut-turut. Dengan satu mutasi tiap sel tiap pembelahan dan satu dari
seratus yang memperbaiki, setidaknya seratus sel harus disaring tiap putaran
untuk menemukan satu; pada praktiknya ribuan, sebab separuh mutasinya memusnahkan
pengikatannya dan seleksinya berisik. Tujuh belas putaran yang masing-masing dua pembelahan
memakan sekitar sembilan hari --- yakni dua sampai tiga minggu yang teramati, mengingat
ketakmangkusannya. Daurnya memisahkan mutasi (zona gelap) dari pengujian
(zona terang) sehingga tiap ragamnya dinilai terhadap \omterm{def:b3:adaptive-immunity:lymphocytes}{antigen} sebelum ia
dibiarkan berlipat ganda, tepat seperti sebuah algoritme genetik menyelang-nyelingkan
ragam dan seleksi.
\end{solution}

\begin{solution}{exo:b3:adaptive-immunity:12}
Perempuan: beberapa gen imun terletak pada kromosom X, sebagian (TLR7)
lolos dari penonaktifan lalu diungkapkan dari kedua salinannya, dan estrogen
meningkatkan ketahanan hidup sel~B --- ujinya: lelaki dengan sebuah X tambahan (Klinefelter)
punya risiko lupus serupa perempuan, dan memang demikian. \omterm{def:b3:adaptive-immunity:mhc}{HLA}: alel
yang berkaitan menyajikan peptida diri yang bersangkutan dengan baik, atau gagal menyajikannya
di dalam \omterm{def:b3:adaptive-immunity:tolerance}{timus} sehingga klonanya tidak dihapus --- ujinya: mencit
\omterm{def:b3:genetic-engineering:transgenic}{transgenik} bagi alel manusianya mengembangkan penyakitnya ketika diberi
\omterm{def:b3:adaptive-immunity:lymphocytes}{antigennya}. Kenaikannya: pemaparan mikroba dini yang lebih sedikit dan
\omterm{def:b3:microbiomes:symbiosis}{mikrobiom} yang berubah (\omterm{def:b3:bacteriology:antibiotics}{antibiotik}, makanan, kelahiran sesar) memberi lebih sedikit sel~T
pengatur dan lebih banyak \omterm{def:b3:innate-immunity:inflammation}{peradangan} --- ujinya: mencit rawan diabetes yang dibesarkan
bebas kuman atau dengan \omterm{def:b3:bacteriology:antibiotics}{antibiotik} lebih banyak mengembangkan diabetes, dan anak
imigran memperoleh angka kejadian negeri barunya dalam satu
angkatan.
\end{solution}

\begin{solution}{pb:b3:adaptive-immunity:1}
\textbf{1.} $6000\times 320 = 1.9\times 10^{6}$; $\times 3\times 10^{4}
= 5.8\times 10^{10}$.
\textbf{2.} $10^{11}$ sel terhadap $5.8\times 10^{10}$ yang mungkin: sebuah tubuh
dapat memuat tiap reseptornya sekitar dua kali, tetapi klonanya tidak sama besar dan
kebanyakan urutannya tidak pernah kebetulan dibuat, sehingga tiap tubuh mengungkapkan
sebuah contoh yang besar tetapi tidak lengkap --- dan dua orang berbagi hanya sebagian
kecil reseptornya.
\textbf{3.} $10^{11}/10^{5} = 10^{6}$ sel; di dalam sebuah kelenjar berisi $10^{8}$,
sekitar $1000$.
\textbf{4.} $10^{9}/10^{5} = 10^{4}$ pendahulu --- yakni seratus kali lebih sedikit;
bayi baru lahir juga punya folikel dan \omterm{def:b3:innate-immunity:innate}{sel dendritik} yang belum matang, tanpa
ingatan, serta pertolongan sel~T yang terbatas, sehingga tanggapannya lambat dan kecil, dan
IgG ibunya menjembatani celahnya.
\textbf{5.} Ia merentang sambungan V--D dan D--J, tempat nukleotida
dibuang dan ditambahkan secara acak: urutannya tidak disandi di mana pun di dalam
genomnya. Karena duduk di antara kedua rantainya di tengah tapak
pengikatnya, ia gelung yang paling sering bersentuhan dengan epitopnya.
\textbf{6.} Sebelum \omterm{def:b3:adaptive-immunity:lymphocytes}{antigennya}, keanekaannya harus luas sehingga sesuatu
mengikat apa pun yang datang; sesudah \omterm{def:b3:adaptive-immunity:lymphocytes}{antigennya}, upayanya lebih baik dihabiskan untuk meragamkan
beberapa reseptor yang sudah diketahui mengikat. Sebuah pencarian kasar lalu yang halus
jauh lebih mangkus daripada salah satunya saja.
\textbf{7.} Tujuh hari adalah $21$ pembelahan: $100\times 2^{21} = 2.1\times
10^{8}$ sel; \qty{30}{\%}, yakni $6.3\times 10^{7}$ \omterm{def:b3:adaptive-immunity:response}{sel plasma}.
\textbf{8.} $6.3\times 10^{7}\times 2000\times \num{86400} = 1.1\times
10^{16}$ molekul sehari; $\times 1.5\times 10^{5}\times 1.66\times
10^{-24}$ g $= \qty{2.7}{mg}$; di dalam \qty{3}{L}, \qty{0.9}{mg/L} tiap hari.
\textbf{9.} \qty{27}{mg} di dalam \qty{3}{L}: \qty{9}{mg/L} $= \qty{0.009}{g/L}$,
yakni seperseribu IgG totalnya --- \omterm{def:b3:adaptive-immunity:antibody}{antibodi} terhadap satu \omterm{def:b3:adaptive-immunity:lymphocytes}{antigen} adalah
sebagian kecil dari keseluruhannya.
\textbf{10.} Seratus kali lipat adalah $\log_{2}100 = 6.6$ paruh hidup, yakni sekitar
$140$ hari. \omterm{def:b3:adaptive-immunity:antibody}{Antibodi} bertahun-tahun datang dari \omterm{def:b3:adaptive-immunity:response}{sel plasma} berumur panjang di dalam
sumsumnya, yang menyekresikan terus-menerus, dan dari \omterm{def:b3:adaptive-immunity:response}{sel ingatan} yang dirangsang ulang
pada pemaparan ulang.
\textbf{11.} Sembilan pembelahan: $10^{5}\times 512 = 5\times 10^{7}$ sel pada
hari 3 --- yakni seribu kali $5\times 10^{4}$ milik tanggapan primer pada hari 3, dan
seperempat dari apa yang dicapai tanggapan primernya baru pada hari 7.
\textbf{12.} VDJ yang tersusun ulang terletak di sebelah C$_{\mu}$, sehingga transkrip
pertamanya IgM; menukar ke C$_{\gamma}$ memerlukan pertolongan sel~T dan AID,
yang memakan waktu berhari-hari. \omterm{def:b3:adaptive-immunity:response}{Sel ingatan} sudah menukar kelasnya, sehingga keturunannya
membuat IgG seketika.
\textbf{13.} $1000\times 10^{-3} = 1$ tiap pembelahan; $20$ sepanjang dua puluh.
\textbf{14.} Satu dari seratus; $10^{4}\times 0.01 = 100$ sel yang membaik
tiap pembelahan.
\textbf{15.} $1.5^{n} = 1000$: $n \approx 17$; $17\times \qty{12}{h}
\approx 8.5$ hari.
\textbf{16.} Separuh anakannya kehilangan pengikatannya: tanpa seleksi, $0.5^{20}
\approx 10^{-6}$ dari garis keturunannya akan masih mengikat sesudah dua puluh
pembelahan. Zona terangnya harus menyingkirkan yang kalah pada tiap daur sehingga
hanya reseptor yang berfungsi dan membaik yang berlanjut.
\textbf{17.} Dosis pertamanya meninggalkan sel~B ingatan berafinitas naik dan
tertukar kelas; dosis keduanya menyemai \omterm{def:b3:adaptive-immunity:germinal-centre}{pusat germinal} baru dengan sel itu,
lalu hipermutasi dan seleksinya berlanjut dari titik awal yang lebih tinggi,
sehingga \omterm{def:b3:adaptive-immunity:antibody}{antibodi} yang muncul lebih baik lagi.
\textbf{18.} Sebagian orang yang terinfeksi akhirnya membuat \omterm{def:b3:adaptive-immunity:antibody}{antibodi} terhadap
tapak lestari yang tidak dapat diubah \omterm{def:b3:virology:virus}{virusnya} (\omterm{def:b3:adaptive-immunity:antibody}{antibodi} penetral
berjangkauan luas), sesudah bertahun-tahun pematangan. \omterm{def:b3:adaptive-immunity:germinal-centre}{Pusat germinalnya} dapat
dikemudikan: sebuah urutan imunogen yang mula-mula melibatkan pendahulu nutfah
yang tepat lalu mengganjar pengikatan pada tapak lestarinya dapat menuntun
pematangannya sepanjang jalan itu dalam hitungan bulan alih-alih tahun.
\textbf{19.} $p_{c} = 1 - 1/15 = 93.3\,\%$. Satu dosis: $0.933/0.93 =
100.4\,\%$ --- tak tercapai; dua dosis: $0.933/0.97 = 96.2\,\%$.
\textbf{20.} Yang kebal $0.90\times 0.97 = 0.873$; $R_{\text{eff}} = 15\times
0.127 = 1.9$: campaknya beredar.
\textbf{21.} $60\,000\times (0.10 + 0.90\times 0.03) = 7600$ orang rentan
tiap angkatan kelahiran; sesudah lima tahun $38\,000$ --- sebuah kumpulan yang di dalamnya sebuah
kasus impor menemukan cukup banyak orang rentan di antara sentuhannya untuk menopang
rantainya.
\textbf{22.} Di antara yang rentan, $s_{\infty} - \ln s_{\infty}/1.9 =
1$ memberi $s_{\infty} \approx 0.23$: sekitar \qty{77}{\%} dari mereka, yakni sekitar
$29\,000$ kasus.
\textbf{23.} Yang kebal $0.95\times 0.97 = 0.92$; $R_{\text{eff}} = 15\times
0.08 = 1.2$ --- masih di atas satu; ambangnya memerlukan \qty{96}{\%} dengan
dua dosis, dan campak adalah penyakit yang menghukum beberapa persen yang
terakhir.
\textbf{24.} Bayi terlindungi hanya oleh imunitas orang di sekitarnya
(dan oleh \omterm{def:b3:adaptive-immunity:antibody}{antibodi} ibunya selama beberapa bulan): bila rantai
penularannya tidak dapat terbentuk, tidak ada kasus yang mencapainya. Pada cakupan \qty{85}{\%}
bagian yang kebal adalah $0.82$ dan $R_{\text{eff}} = 2.6$: letusannya
berjalan, dan bayinya, yang tidak tervaksinasi dan paling rentan, terinfeksi.
\textbf{25.} Khazanahnya sekitar $6\times 10^{10}$; \qty{2.7}{mg} IgG
sehari dari \omterm{def:b3:adaptive-immunity:response}{sel plasma} satu tanggapan; cakupan dua dosis \qty{96}{\%}.
\end{solution}
